\documentclass[11pt]{article}
\usepackage[margin=1in]{geometry}
%%%%%%%%%%%%%%%%%%%%%%%%%%%%%%%%%%%%%%%%%%%%%%%%%%%%%%%%%%%%%%%%%%%%%%%%%%%%
\usepackage{amsmath,amssymb,amsfonts}
\usepackage{amsthm,latexsym,mathrsfs}
\usepackage{bm,bbm,dsfont}
\usepackage{color,graphicx}
\usepackage[latin1]{inputenc}
\usepackage[english]{babel}
\usepackage{algorithm}
\usepackage{algpseudocode}
\usepackage{pgfplots}
\pgfplotsset{compat=1.18}
\usepgfplotslibrary{groupplots}
% the colors of the locations, as in the main document
\definecolor{locNE}{HTML}{2A78D6}
\definecolor{locMW}{HTML}{EB6834}
\definecolor{locS}{HTML}{1BAF7A}
\definecolor{locW}{HTML}{EDA100}
\usepackage{xr-hyper}
\usepackage[colorlinks=true,linkcolor=blue,citecolor=blue,urlcolor=blue]{hyperref}
% labels of the main document, prefixed M-: build it first, so that its .aux exists
\externaldocument[M-]{math_bio_positive_matrices_clean}

\theoremstyle{plain}
\newtheorem{theorem}{Theorem}[section]
\newtheorem{proposition}[theorem]{Proposition}
\theoremstyle{remark}
\newtheorem{remark}[theorem]{Remark}

% section numbers in headings followed by a period
\makeatletter
\renewcommand\@seccntformat[1]{\csname the#1\endcsname.\quad}
\makeatother

\everymath{\displaystyle}
\numberwithin{equation}{section}
\renewcommand{\thesection}{T\arabic{section}}

\title{Three doses: block Gauss--Seidel and the joint method\thanks{Parts of this manuscript were edited with the assistance of an AI tool (Claude Opus 5.5, Anthropic). The authors are solely responsible for all proofs and statements.}}
\author{Zhilan \and John \and Yuliya \and Ludmil}
\date{}

\begin{document}

\maketitle

\begin{abstract}
Companion to ``Minimizing the spectral radius of a diagonally scaled positive
matrix''. Equation, theorem and section numbers without a prefix T refer to
that document. For the three-dose problem~\eqref{M-three-program} we study the
minimization over one dose type at a time with the other two fixed, in
Gauss--Seidel order, and compare it with a local method in all variables, from
the barycenter of the feasible set and from $48$ starts near its vertices.
\end{abstract}

\section{The problem and its blocks}\label{t-blocks}

We consider~\eqref{M-three-program} with $b_{k,i}>0$ for all $k=1:3$ and
$i=1:m$. By property~(\ref{M-simplex-product}) of Section~\ref{M-sec-reduced}
the feasible set is
$\mathcal{S}=\mathcal{S}_1\times\mathcal{S}_2\times\mathcal{S}_3$, with
$\mathcal{S}_k=\left(\Delta_{n-1}\right)^m$ the set of admissible
$\widehat{\bm u}_k=\left(\widehat u_{k,ij}\right)$, and the objective
function is
\begin{equation}\label{t-F}
F(\widehat{\bm u}_1,\widehat{\bm u}_2,\widehat{\bm u}_3)
=\log\rho\left(A_1\operatorname{diag}\left(\bm g\left(\bm\chi\right)\right)\right),
\qquad \chi_{k,p}=\frac{b_{k,i}}{\sigma_{k,ij}}\,\widehat u_{k,ij},
\end{equation}
with $\bm g$ given by~\eqref{M-g-def}. With the eigenvectors $\bm v,\bm w>0$
of~\eqref{M-perron}, the argument that gives~\eqref{M-grad} gives
\begin{equation}\label{t-grad}
\frac{\partial F}{\partial\widehat u_{k,ij}}
=\frac{1}{\rho}\cdot\frac{\left(A_1^T\bm w\right)_pv_p}{\langle\bm w,\bm v\rangle}
\cdot\frac{\partial g_p}{\partial \chi_{k,p}}\cdot\frac{b_{k,i}}{\sigma_{k,ij}},
\qquad p=(i-1)n+j .
\end{equation}
Minimizing $F$ over one block $\widehat{\bm u}_k\in\mathcal{S}_k$ with the
other two fixed is the step of the method of Section~\ref{t-algorithm}. The
three blocks behave differently.

\begin{proposition}\label{t-block1}
Let $A_1\ge 0$ be irreducible and let $\widehat{\bm u}_2$ and
$\widehat{\bm u}_3$ be fixed. Then
$\widehat{\bm u}_1\mapsto F(\widehat{\bm u}_1,\widehat{\bm u}_2,\widehat{\bm u}_3)$
is strictly convex on $\mathcal{S}_1$ and has a unique minimizer there.
\end{proposition}
\begin{proof}
With $\chi_{2,p}$ and $\chi_{3,p}$ fixed, $\pi_{2,p}$ and $\pi_{3,p}$ in~\eqref{M-g-def}
are constants, and $\pi_{1,p}=\chi_{1,p}/(\omega_v+\varepsilon_2\chi_{2,p})$ is linear
in $\chi_{1,p}$. Put $a_p=(1-\varepsilon_1)+\pi_{2,p}$ and
$\beta_p=1+\pi_{2,p}+\pi_{2,p}\pi_{3,p}$. Then
\begin{equation}
g_p=\frac{1+a_p\pi_{1,p}}{1+\beta_p\pi_{1,p}}=c_p+\frac{d_p}{\tau_p},\qquad
c_p=\frac{a_p}{\beta_p},\quad d_p=1-\frac{a_p}{\beta_p},\quad
\tau_p=1+\beta_p\pi_{1,p}.
\end{equation}
Here $a_p\ge 1-\varepsilon_1\ge 0$ and $\beta_p-a_p=\varepsilon_1+\pi_{2,p}\pi_{3,p}>0$,
so $c_p\ge 0$ and $d_p>0$, and $\tau_p\ge 1$ is affine in $\widehat u_{1,ij}$
with positive slope. This is the form $g_p=(1-\varepsilon_1)+\varepsilon_1/t_p$
of the single dose, with constants depending on $p$, and the proof of
Theorem~\ref{M-global} applies without change: $\log(c_p+d_p/\tau_p)$ is
strictly convex in $\tau_p$ for $c_p\ge 0$, $d_p>0$; the function $\Phi$ is convex
by Lemma~\ref{M-log-convex} in~\S\ref{M-sec-facts} and strictly increasing in each argument by
property~(\ref{M-mono-g}) of Section~\ref{M-sec-reduced}. Hence $F$ is strictly convex in $\widehat{\bm u}_1$, and
on the compact convex set $\mathcal{S}_1$ it has a unique minimizer.
\end{proof}

\begin{proposition}\label{t-block3}
Let $\widehat{\bm u}_1$ and $\widehat{\bm u}_2$ be fixed. Then
$\widehat{\bm u}_3\mapsto F(\widehat{\bm u}_1,\widehat{\bm u}_2,\widehat{\bm u}_3)$
is convex on $\mathcal{S}_3$.
\end{proposition}
\begin{proof}
The numerator of $g_p$ in~\eqref{M-g-def} does not depend on $\chi_{3,p}$, and the
denominator equals $\alpha_p+\gamma_p\chi_{3,p}$ with
$\alpha_p=1+\pi_{1,p}+2\pi_{1,p}\pi_{2,p}>0$ and
$\gamma_p=\varepsilon_3\pi_{1,p}\pi_{2,p}/\omega\ge 0$. So
$\log g_p=\mathrm{const}-\log\left(\alpha_p+\gamma_p\chi_{3,p}\right)$ is convex in
$\chi_{3,p}$, hence in $\widehat u_{3,ij}$. The function $\Phi$ is convex
(see Lemma~\ref{M-log-convex} in~\S\ref{M-sec-facts}) and non-decreasing in
each argument (see Theorem~\ref{M-mono} in~\S\ref{M-sec-facts}), so $F=\Phi\circ\log\bm g$ is convex in
$\widehat{\bm u}_3$.
\end{proof}

If $\gamma_p=0$, that is, if $\chi_{1,p}=0$ or $\chi_{2,p}=0$, then $g_p$ does not
depend on $\chi_{3,p}$, and the minimizer over $\mathcal{S}_3$ need not be unique.

For the second block neither convexity nor monotonicity holds. The dose $\chi_{2,p}$
enters both $\pi_{1,p}$, through $\omega_v+\varepsilon_2\chi_{2,p}$, and
$\pi_{2,p}$. For $200\,000$ points $(\chi_{1,p},\chi_{2,p},\chi_{3,p})$ drawn
log-uniformly from $\left[10^{-5},10^{-1.5}\right]^3$, with the parameters of
the model, $\varepsilon_1=\varepsilon_2=\varepsilon_3=0.8$ and
$\omega=\omega_v=1/182.5$, the smallest value of
$\chi_{2,p}^2\,\partial^2\log g_p/\partial \chi_{2,p}^2$ is $-0.154$, and $g_p$
increases with $\chi_{2,p}$ at $24.0$ per cent of the points. The minimization
over $\mathcal{S}_2$ is therefore a non-convex problem, and its equality
constraints cannot be relaxed to inequalities as in Section~\ref{M-sec-single}.

\section{The algorithms}\label{t-algorithm}

Each minimization, over one block or over all three, is carried out by the
trust region interior point method of~\cite{1999Byrd_Hribar_Nocedal}, with the
gradient~\eqref{t-grad}, a quasi-Newton Hessian, and the equality constraints
$\sum_j\widehat u_{k,ij}=1$; the derivatives $\partial g_p/\partial \chi_{k,p}$
in~\eqref{t-grad} are computed by the complex step. The interior point method
leaves entries of order $10^{-9}$ below zero; after each block minimization
these are set to zero and each $\widehat{\bm u}_{k,i\cdot}$ is rescaled to sum
to one.

\begin{algorithm}[ht]
\caption{Block Gauss--Seidel over the dose types}\label{t-alg-gs}
\begin{algorithmic}[1]
\Require a start $\widehat{\bm u}^{(0)}\in\mathcal{S}$, a tolerance
$\mathrm{tol}$, a maximal number of sweeps $s_{\max}$
\State $\rho_0\gets\exp F\left(\widehat{\bm u}^{(0)}\right)$
\For{$s=1,\ldots,s_{\max}$}
  \For{$k=1,2,3$}
    \State $\widehat{\bm u}_k\gets\operatorname*{arg\,min}_{\bm\eta\in\mathcal{S}_k}
    F(\ldots,\bm\eta,\ldots)$, the other two blocks fixed, started from the current $\widehat{\bm u}_k$
    \State set the negative entries of $\widehat{\bm u}_k$ to zero and rescale each simplex to sum one
  \EndFor
  \State $\rho_s\gets\exp F\left(\widehat{\bm u}\right)$
  \If{$\left|\rho_s-\rho_{s-1}\right|<\mathrm{tol}\cdot\rho_{s-1}$}
    \State \textbf{break}
  \EndIf
\EndFor
\State verify the Karush--Kuhn--Tucker conditions at $\widehat{\bm u}$ (Algorithm~\ref{t-alg-kkt})
\State \Return $\widehat{\bm u}$, $\rho_s$
\end{algorithmic}
\end{algorithm}

\begin{algorithm}[ht]
\caption{Verification of the Karush--Kuhn--Tucker conditions}\label{t-alg-kkt}
\begin{algorithmic}[1]
\Require a computed point $\widehat{\bm u}\in\mathcal{S}$, a threshold $\theta$
\State compute $\nabla F\left(\widehat{\bm u}\right)$ by~\eqref{t-grad}
\For{each factor $(k,i)$}
  \State $J_+\gets\{j:\widehat u_{k,ij}>\theta\}$, $J_0\gets\{j:\widehat u_{k,ij}\le\theta\}$
  \State $\lambda_{k,i}\gets$ mean of $\partial F/\partial\widehat u_{k,ij}$ over $j\in J_+$
  \State $\mathrm{spread}_{k,i}\gets\left(\max_{J_+}-\min_{J_+}\right)\partial F/\partial\widehat u_{k,ij}\,/\,|\lambda_{k,i}|$
  \State $\mathrm{excess}_{k,i}\gets\min_{j\in J_0}\left(\partial F/\partial\widehat u_{k,ij}-\lambda_{k,i}\right)/|\lambda_{k,i}|$
\EndFor
\State \Return $\max_{k,i}\mathrm{spread}_{k,i}$ and $\min_{k,i}\mathrm{excess}_{k,i}$
\end{algorithmic}
\end{algorithm}

At a point satisfying the Karush--Kuhn--Tucker conditions the spread is zero
and the excess is non-negative.

The \emph{joint method}, Algorithm~\ref{t-alg-joint}, minimizes $F$ in all $3N$
variables $\left(\widehat{\bm u}_1,\widehat{\bm u}_2,\widehat{\bm u}_3\right)$
at once, with the same interior point method and the gradient~\eqref{t-grad} in
all three blocks; Algorithm~\ref{t-alg-gs} uses the same method in $N$
variables at a time.

\begin{algorithm}[ht]
\caption{The joint method}\label{t-alg-joint}
\begin{algorithmic}[1]
\Require a start $\widehat{\bm u}^{(0)}\in\mathcal{S}$
\State $\widehat{\bm u}\gets\operatorname*{arg\,min}_{\bm\eta\in\mathcal{S}}F(\bm\eta)$,
started from $\widehat{\bm u}^{(0)}$, with the constraints
$\sum_j\eta_{k,ij}=1$ and $\bm\eta\ge 0$ for all $k=1:3$, $i=1:m$
\State set the negative entries of $\widehat{\bm u}$ to zero and rescale each simplex to sum one
\State verify the Karush--Kuhn--Tucker conditions at $\widehat{\bm u}$ (Algorithm~\ref{t-alg-kkt})
\State \Return $\widehat{\bm u}$, $\rho=\exp F\left(\widehat{\bm u}\right)$
\end{algorithmic}
\end{algorithm}

Each block step of Algorithm~\ref{t-alg-gs} does not increase $F$ when the
block minimization is exact, so $\rho_s$ does not increase. If, in addition,
every block minimization has a unique solution, every limit point of the
iterates is a stationary point of~\eqref{M-three-program};
see~\cite[Section 2.7]{1999Bertsekas}. Proposition~\ref{t-block1} gives this
uniqueness for the first block; for the second and the third it is not
guaranteed.

\begin{algorithm}[ht]
\caption{The comparison over starts}\label{t-alg-starts}
\begin{algorithmic}[1]
\Require $\delta\in(0,1)$
\State $\overline{\bm u}\gets$ the barycenter of $\mathcal{S}$, $\overline u_{k,ij}=1/n$
\State starts $\gets\{\overline{\bm u}\}$
\For{$k=1:3$, $i=1:m$, $j=1:n$}
  \State $\widehat{\bm u}\gets\overline{\bm u}$ with the factor $(k,i)$ replaced by
  $(1-\delta)\bm e_j+\delta\,\mathbbm{1}/n$
  \State add $\widehat{\bm u}$ to the starts
\EndFor
\For{each start}
  \State run Algorithm~\ref{t-alg-joint} and Algorithm~\ref{t-alg-gs}; record $\rho$, the
  Karush--Kuhn--Tucker check, iterations or sweeps, and the time
\EndFor
\end{algorithmic}
\end{algorithm}

\section{Results}\label{t-results}

The instance is that of Section~S3 of the supplement with the weekly budgets:
$b_{1i}$ the weekly budget of the first dose, $b_{2i}=b_{3i}=b_{1i}/4$, and
$\sigma_{2,ij}=\sigma_{3,ij}$ the numbers of partially protected individuals,
no budget for the second and third doses being given with the data
of~\cite{JZh}. The parameters are $m=n=4$, $N=16$, $\delta=10^{-3}$,
$\mathrm{tol}=10^{-12}$, $s_{\max}=200$ and $\theta=10^{-8}$, so that
Algorithm~\ref{t-alg-starts} has $1+3mn=49$ starts. Table~\ref{t-tab} collects
the results.

\begin{table}[ht]
\centering
\small
\begin{tabular}{lllll}
\hline
 & joint method & \multicolumn{2}{l}{Algorithm~\ref{t-alg-gs}}\\
starts & all $49$ & barycenter, and dose $1$ or $2$ & dose $3$ at a vertex\\
 & & at a vertex ($1+32$) & ($16$)\\
\hline
$\rho$ & $1.6618260420$--$1.6618260490$ & $1.6618260420$ & $16$ values, $1.7727$--$1.9510$\\
spread & not conclusive & $\le 6.6\cdot 10^{-7}$ & $\le 1.0\cdot 10^{-6}$ at $15$; $4.0$ at one\\
excess & not conclusive & $\ge 1.0\cdot 10^{-3}$ & $\ge 5.5\cdot 10^{-2}$ at the $15$\\
iterations & $128$--$1924$ & $9$--$37$ sweeps & $9$--$79$ sweeps\\
time & $25$ s & $165$ s & $163$ s\\
\hline
\end{tabular}
\caption{The joint method and Algorithm~\ref{t-alg-gs} from the $49$ starts of
Algorithm~\ref{t-alg-starts}. A start with ``dose $k$ at a vertex'' is the
barycenter with one factor $(k,i)$ replaced by $(1-\delta)\bm e_j+\delta\mathbbm{1}/n$:
in location $i$ all doses of type $k$ go, up to $\delta$, to the age group $j$.
There are $mn=16$ such starts for each $k$. The rows ``spread'' and ``excess'' are the output
of Algorithm~\ref{t-alg-kkt}; the times are totals over the starts of the
column.}\label{t-tab}
\end{table}

The joint method gives the same value, to $7\cdot 10^{-9}$, from every start.
Its endpoints keep a few entries between $10^{-8}$ and $5\cdot 10^{-6}$ where
the endpoints of Algorithm~\ref{t-alg-gs} have zeros: from the barycenter and
the starts $(k,i,j)=(1,1,4)$ and $(2,1,1)$ there are $1$, $6$ and $2$ such
entries, and the two endpoints differ by at most $4.6\cdot 10^{-6}$.
Algorithm~\ref{t-alg-kkt} with $\theta=10^{-8}$ counts these entries as
positive, and its spread is then not small, so for the joint method the check
is not conclusive. The smallest value,
$\rho=1.6618260420$, is attained by Algorithm~\ref{t-alg-gs} from the
barycenter and from all $32$ starts with dose $1$ or $2$ at a vertex, at points
that pass the check.

From the $16$ starts with dose $3$ at a vertex, in which the third dose of one location is
concentrated in one age group, Algorithm~\ref{t-alg-gs} ends at $16$ different
values between $1.7727$ and $1.9510$, $15$ of them at points that pass the
check; in the remaining one the spread is $4.0$ after $37$ sweeps. Two of
these endpoints, from the starts $(k,i,j)=(3,1,4)$ and $(3,2,1)$, with
$\rho=1.9509854917$ and $\rho=1.8308005079$, were examined further. For each,
none of $2000$ random feasible directions decreases $F$, at steps $10^{-5}$,
$10^{-4}$, $10^{-3}$ and $10^{-2}$, so that these are local minimizers to the
resolution of the test; the joint method started at the endpoint, moved
$10^{-6}$ towards the barycenter, reaches $\rho=1.6618260434$. At the first of
them, in location $1$ all of the second and third doses are given to the age
group $60+$; there $\chi_{2,p}=0$ in the other three groups of location $1$, so
$\gamma_p=0$ in Proposition~\ref{t-block3} and the derivative of $F$ in the
third dose vanishes in these groups, while the derivative in the second dose
exceeds the multiplier. A step in one block alone does not decrease $F$; the
decrease to the smaller value requires the second and third doses to move
together.

The same endpoint shows that $F$ is not convex on this instance: at the
midpoint of the segment from it to the minimizer of the joint method,
$\log\rho=0.590576$, against $0.588126$ for the mean of the values at the two
ends. The random-segment test of Section~S3 of the supplement did not find a
violation.

\section{Is the barycenter a reliable start?}\label{t-instances}

To test whether the barycenter of $\mathcal{S}$ is a reliable start, we varied
the data in $36$ instances: the first-dose budget $b_{1i}$ weekly or daily
($b_{1i}/7$); the ratios
\begin{equation}
\frac{(b_{2i},b_{3i})}{b_{1i}}\in\left\{\left(\tfrac1{16},\tfrac1{16}\right),\left(\tfrac14,\tfrac14\right),(1,1),(4,4),\left(\tfrac12,\tfrac18\right),\left(\tfrac18,\tfrac12\right)\right\};
\end{equation}
and $\sigma_{2,ij}$, $\sigma_{3,ij}$ equal to the numbers $P_{ij}$ of partially
protected individuals, or equal to $P_{ij}e^{Z_{ij}/2}$ and $P_{ij}e^{Z'_{ij}/2}$
with independent standard normal $Z$, $Z'$, for two seeds. For each instance
the reference value $\rho_{\mathrm{ref}}$ is the smallest $\rho$ of
Algorithm~\ref{t-alg-joint} over $69$ starts: the $49$ of
Algorithm~\ref{t-alg-starts} and $20$ random points of $\mathcal{S}$. From the
barycenter we ran Algorithm~\ref{t-alg-joint} and Algorithm~\ref{t-alg-gs}, and
call a run successful if its $\rho$ is at most $(1+10^{-7})\rho_{\mathrm{ref}}$.
Table~\ref{t-tab-inst} summarizes the results, and Table~\ref{t-tab-fail} lists
the instances in which a run from the barycenter was not successful.

\begin{table}[ht]
\centering
\small
\begin{tabular}{lcc}
\hline
 & weekly $b_{1i}$ & daily $b_{1i}/7$\\
\hline
instances & $18$ & $18$\\
instances with more than one value among the $69$ starts & $2$ & $12$\\
largest number of distinct values & $3$ & $4$\\
Algorithm~\ref{t-alg-joint} from the barycenter successful & $17$ & $16$\\
Algorithm~\ref{t-alg-gs} from the barycenter successful & $16$ & $12$\\
at least one of the two successful & $17$ & $17$\\
\hline
\end{tabular}
\caption{Runs from the barycenter on the $36$ instances. Values are distinct if
they differ by more than $10^{-6}$ relative.}\label{t-tab-inst}
\end{table}

\begin{table}[ht]
\centering
\small
\begin{tabular}{llccccc}
\hline
budget & $(b_{2i},b_{3i})/b_{1i}$ & $\sigma_2,\sigma_3$ &
Alg.~\ref{t-alg-joint} & Alg.~\ref{t-alg-gs} & distinct & starts at $\rho_{\mathrm{ref}}$\\
\hline
weekly & $(\tfrac1{16},\tfrac1{16})$ & $P$       & $3\cdot 10^{-14}$ & $1.2\cdot 10^{-2}$ & $3$ & $93\%$\\
weekly & $(\tfrac1{16},\tfrac1{16})$ & seed $2$  & $5.4\cdot 10^{-3}$ & $5.4\cdot 10^{-3}$ & $2$ & $1\%$\\
daily  & $(\tfrac1{16},\tfrac1{16})$ & $P$       & $1.1\cdot 10^{-4}$ & $1.1\cdot 10^{-4}$ & $3$ & $28\%$\\
daily  & $(\tfrac14,\tfrac14)$       & seed $1$  & $3\cdot 10^{-12}$ & $5.0\cdot 10^{-3}$ & $2$ & $99\%$\\
daily  & $(\tfrac12,\tfrac18)$       & $P$       & $7\cdot 10^{-11}$ & $1.2\cdot 10^{-7}$ & $2$ & $99\%$\\
daily  & $(\tfrac12,\tfrac18)$       & seed $1$  & $3\cdot 10^{-12}$ & $2.2\cdot 10^{-3}$ & $2$ & $91\%$\\
daily  & $(\tfrac12,\tfrac18)$       & seed $2$  & $3\cdot 10^{-12}$ & $1.3\cdot 10^{-7}$ & $1$ & $100\%$\\
daily  & $(\tfrac18,\tfrac12)$       & $P$       & $2.2\cdot 10^{-3}$ & $8\cdot 10^{-8}$ & $3$ & $1\%$\\
daily  & $(\tfrac18,\tfrac12)$       & seed $1$  & $3\cdot 10^{-12}$ & $8.4\cdot 10^{-3}$ & $2$ & $99\%$\\
\hline
\end{tabular}
\caption{The instances in which a run from the barycenter was not successful:
the relative excess $(\rho-\rho_{\mathrm{ref}})/\rho_{\mathrm{ref}}$ of
Algorithms~\ref{t-alg-joint} and~\ref{t-alg-gs}, the number of distinct values
among the $69$ starts, and the fraction of the starts that reach
$\rho_{\mathrm{ref}}$.}\label{t-tab-fail}
\end{table}

The barycenter is not always a good start. Algorithm~\ref{t-alg-joint} from it
misses $\rho_{\mathrm{ref}}$ in $3$ of the $36$ instances, by up to $0.54$ per
cent, and Algorithm~\ref{t-alg-gs} in $8$, by up to $1.2$ per cent; two of the
latter, with excess $1.2\cdot 10^{-7}$ and $1.3\cdot 10^{-7}$, are a matter of
accuracy rather than of a different minimizer. Taking the better of the two
runs from the barycenter misses $\rho_{\mathrm{ref}}$ in $2$ instances, both
with the smallest budgets $b_{2i}=b_{3i}=b_{1i}/16$. Where Algorithm~\ref{t-alg-joint} fails from the
barycenter, only $1$ to $28$ per cent of the $69$ starts reach
$\rho_{\mathrm{ref}}$, so the smallest value is attained from a small part of
$\mathcal{S}$. The problem is also less benign for the daily budgets: $12$ of
the $18$ daily instances have more than one value among the starts, against
$2$ of the $18$ weekly ones.

In $11$ of the daily instances Algorithm~\ref{t-alg-kkt} reports a spread above
$3$ at the endpoint of Algorithm~\ref{t-alg-gs}, also where this endpoint is
at $\rho_{\mathrm{ref}}$; as for the joint method in Section~\ref{t-results},
the check with $\theta=10^{-8}$ is then not conclusive. The value
$\rho_{\mathrm{ref}}$ itself is the best of $69$ local runs, not a certified
global minimum. Two instances took much longer than the others,
$2336$~s and $1329$~s against $10$ to $60$~s, because of some of the $69$ runs
of Algorithm~\ref{t-alg-joint}.

\section{An adaptive first-order method}\label{t-yogi}

The joint method and Algorithm~\ref{t-alg-gs} use a quasi-Newton Hessian. We
also tested a first-order method in all $3N$ variables: the adaptive method
Yogi of~\cite{2018Zaheer_Yogi}, a variant of Adam~\cite{2015Kingma_Ba}, with
each step projected onto $\mathcal{S}$ in the metric of its diagonal
preconditioner. The method uses the gradient~\eqref{t-grad} and no Hessian.

\subsection{Smoothness of the objective function}

For $\widehat{\bm u}\in\mathcal{S}$ all $\chi_{k,p}\ge 0$, so
$\pi_{1,p},\pi_{2,p}\ge 0$ and $\pi_{3,p}\ge 1$ in~\eqref{M-g-def}; the
numerator of $g_p$ is at least $1$ and its denominator at least $1$, and $g_p$
is a rational function of $\bm\chi$ with a positive denominator on a
neighbourhood of $\mathcal{S}$. Hence $g_p>0$ there, and
$A_1\operatorname{diag}(\bm g)$ is non-negative and irreducible with the
pattern of $A_1$. Its spectral radius is a simple eigenvalue, so it is an
analytic function of the entries near every such $\bm g$
(see~\cite[Section 6.3]{1985Horn_Johnson}), and $F$ is $C^\infty$ on a
neighbourhood of the compact set $\mathcal{S}$. In particular $\nabla F$,
given by~\eqref{t-grad}, is continuous and bounded on $\mathcal{S}$, and we
put
\begin{equation}\label{t-G}
G=\max_{\widehat{\bm u}\in\mathcal{S}}\left\|\nabla F\left(\widehat{\bm u}\right)\right\|_\infty<\infty .
\end{equation}

\subsection{Projection in a diagonal metric}

For a vector $\bm d\in\mathbb{R}^n$ with $\bm d>0$ and $\bm z\in\mathbb{R}^n$ let
\begin{equation}\label{t-proj}
\Pi^{\bm d}(\bm z)=\operatorname*{arg\,min}_{\bm y\in\Delta_{n-1}}\
\frac12\sum_{j=1}^n d_j\left(y_j-z_j\right)^2 .
\end{equation}
For $\mathcal{S}$, a product of $3m$ simplices, and a diagonal matrix $D>0$
of order $3N$, $\Pi^D_{\mathcal{S}}$ applies~\eqref{t-proj} to each factor
$(k,i)$ with the diagonal entries of $D$ of that factor. With $D=I$ this is
the Euclidean projection $\Pi_{\mathcal{S}}$.

\begin{proposition}\label{t-proj-prop}
Let $\bm d>0$ and $\bm z\in\mathbb{R}^n$.
\begin{enumerate}
\item[(a)] Problem~\eqref{t-proj} has a unique solution $\bm y$, and
$\bm y\in\Delta_{n-1}$ is this solution if and only if
\begin{equation}\label{t-vi}
\sum_{j=1}^n d_j\left(z_j-y_j\right)\left(x_j-y_j\right)\le 0
\quad\text{for all }\bm x\in\Delta_{n-1}.
\end{equation}
\item[(b)] The solution is $y_j=\max\left(z_j-\lambda/d_j,0\right)$, where
$\lambda$ is the unique root of
\begin{equation*}
\varphi(\lambda)=\sum_j\max\left(z_j-\lambda/d_j,0\right)=1 .
\end{equation*}
\item[(c)] $\Pi^{\bm d}(\bm z)$ is a continuous function of $(\bm z,\bm d)$
on $\mathbb{R}^n\times(0,\infty)^n$.
\end{enumerate}
\end{proposition}
\begin{proof}
(a) The objective function of~\eqref{t-proj} is continuous and strictly convex,
and $\Delta_{n-1}$ is compact and convex, so the minimizer exists and is
unique. For a convex differentiable function on a convex set, a point is a
minimizer if and only if the directional derivative towards every point of the
set is non-negative~\cite[Section 4.2.3]{2004Boyd_Vandenberghe}; here this
derivative is $\sum_j d_j(y_j-z_j)(x_j-y_j)$, which is~\eqref{t-vi}.

(b) The function $\varphi$ is continuous and non-increasing, it is
$0$ for $\lambda\ge\max_j d_jz_j$, and $\varphi(\lambda)\to\infty$ as
$\lambda\to-\infty$. Where $\varphi(\lambda)>0$ at least one term is positive,
and that term is strictly decreasing in $\lambda$, so $\varphi$ is strictly
decreasing on $\{\varphi>0\}$ and $\varphi(\lambda)=1$ has exactly one root.
For this $\lambda$ the vector $\bm y$ lies in $\Delta_{n-1}$. For
$\bm x\in\Delta_{n-1}$, $d_j(z_j-y_j)=\lambda$ if $y_j>0$, and
$d_j(z_j-y_j)=d_jz_j\le\lambda$ if $y_j=0$; hence
\begin{equation}
\sum_j d_j(z_j-y_j)(x_j-y_j)
=\lambda\sum_{y_j>0}(x_j-y_j)+\sum_{y_j=0}d_jz_jx_j
\le\lambda\sum_{y_j>0}(x_j-y_j)+\lambda\sum_{y_j=0}x_j
=\lambda\left(1-1\right)=0,
\end{equation}
using $x_j\ge 0$, $\sum_jx_j=1$ and $\sum_{y_j>0}y_j=1$. By (a), $\bm y$ is the
solution.

(c) The function $\varphi(\lambda;\bm z,\bm d)$ is jointly continuous. Let
$(\bm z_\ell,\bm d_\ell)\to(\bm z,\bm d)$, with roots $\lambda_\ell$ and
$\lambda$. For $\delta>0$, $\varphi(\lambda-\delta;\bm z,\bm d)>1>\varphi(\lambda+\delta;\bm z,\bm d)$
by the strict decrease in (b), so for $\ell$ large the same inequalities hold
with $(\bm z_\ell,\bm d_\ell)$, and $|\lambda_\ell-\lambda|<\delta$. Hence
$\lambda_\ell\to\lambda$, and the formula in (b) is continuous in
$(\lambda,\bm z,\bm d)$.
\end{proof}

The root in (b) is computed exactly: on each interval between consecutive
values of $d_jz_j$, $\varphi$ is affine in $\lambda$.

\begin{proposition}\label{t-fixed}
Let $D>0$ be diagonal of order $3N$, $\eta>0$, and
$\widehat{\bm u}\in\mathcal{S}$. The following are equivalent:
\begin{enumerate}
\item[(a)] $\widehat{\bm u}=\Pi^D_{\mathcal{S}}\left(\widehat{\bm u}-\eta D^{-1}\nabla F\left(\widehat{\bm u}\right)\right)$;
\item[(b)] $\nabla F\left(\widehat{\bm u}\right)^T\left(\bm x-\widehat{\bm u}\right)\ge 0$ for all $\bm x\in\mathcal{S}$;
\item[(c)] for each factor $(k,i)$ there is $\lambda_{k,i}$ with
$\partial F/\partial\widehat u_{k,ij}=\lambda_{k,i}$ where $\widehat u_{k,ij}>0$
and $\partial F/\partial\widehat u_{k,ij}\ge\lambda_{k,i}$ where $\widehat u_{k,ij}=0$,
\end{enumerate}
that is, $\widehat{\bm u}$ satisfies the Karush--Kuhn--Tucker conditions
of~\eqref{M-three-program}. In particular the condition does not depend on $D$
and $\eta$.
\end{proposition}
\begin{proof}
(a)$\Leftrightarrow$(b). Apply Proposition~\ref{t-proj-prop}(a) to each factor
with
\begin{equation*}
\bm y=\widehat{\bm u}_{k,i\cdot},\qquad
\bm z=\widehat{\bm u}_{k,i\cdot}-\eta D_{k,i}^{-1}\nabla_{k,i}F,
\end{equation*}
where the subscript $k,i$ denotes the entries of the factor. Then
$d_j(z_j-y_j)=-\eta\,\partial F/\partial\widehat u_{k,ij}$, and~\eqref{t-vi}
holding in every factor is the statement that
$\nabla_{k,i}F^T(\bm x_{k,i}-\widehat{\bm u}_{k,i\cdot})\ge 0$ for all
$\bm x_{k,i}\in\Delta_{n-1}$ and all $(k,i)$. Since $\mathcal{S}$ is the
product of the factors, this is (b).

(b)$\Leftrightarrow$(c). Condition (b) says that $\widehat{\bm u}_{k,i\cdot}$
minimizes the linear function $\bm c^T\bm x$, $\bm c=\nabla_{k,i}F$, over
$\Delta_{n-1}$ in every factor. The minimum of $\bm c^T\bm x$ over
$\Delta_{n-1}$ is $\min_jc_j$, and
$\bm c^T\bm y-\min_jc_j=\sum_jy_j\left(c_j-\min_lc_l\right)$, which vanishes for
$\bm y\in\Delta_{n-1}$ if and only if $y_j=0$ whenever $c_j>\min_lc_l$. This
is (c) with $\lambda_{k,i}=\min_jc_j$. Conversely, under (c),
$\lambda_{k,i}\le c_j$ for all $j$ with equality on the support of $\bm y$,
which has at least one element, so $\lambda_{k,i}=\min_jc_j$.

Finally, the Karush--Kuhn--Tucker conditions of~\eqref{M-three-program} in the
variables $\widehat{\bm u}$, with the constraints $\sum_j\widehat u_{k,ij}=1$
and $\widehat u_{k,ij}\ge 0$, are (c) with the multiplier $\lambda_{k,i}$ of
the equality constraint of the factor.
\end{proof}

With $D=I$ and $\eta=1$ we use
\begin{equation}\label{t-r}
r\left(\widehat{\bm u}\right)=\left\|\widehat{\bm u}-\Pi_{\mathcal{S}}\left(\widehat{\bm u}-\nabla F\left(\widehat{\bm u}\right)\right)\right\|_\infty
\end{equation}
as the measure of stationarity: $r\ge 0$ on $\mathcal{S}$, it is continuous,
and by Proposition~\ref{t-fixed} it vanishes exactly at the points satisfying
the Karush--Kuhn--Tucker conditions. Unlike Algorithm~\ref{t-alg-kkt} it does
not need a threshold $\theta$ to decide which entries are positive.

\subsection{The method}

Algorithm~\ref{t-alg-yogi} states the method. Squares, square roots, products,
quotients and $\operatorname{sign}$ of vectors are taken entrywise. Compared
with the method as stated in~\cite{2018Zaheer_Yogi}, the gradient is exact
rather than stochastic, and the step is followed by the projection; we start
the second moment $\bm s$ at the square of the first gradient, and divide the
first moment by $1-\beta_1^t$ as in~\cite{2015Kingma_Ba}.

\begin{algorithm}[ht]
\caption{Projected Yogi}\label{t-alg-yogi}
\begin{algorithmic}[1]
\Require a start $\widehat{\bm u}^{(0)}\in\mathcal{S}$; a step $\eta>0$;
$\beta_1\in[0,1)$, $\beta_2\in(0,1)$, $\epsilon>0$; a tolerance $\mathrm{tol}$
and a maximal number of iterations $t_{\max}$
\State $\bm q\gets\nabla F\left(\widehat{\bm u}^{(0)}\right)$, $\bm \mu\gets\bm 0$, $\bm s\gets\bm q^2$
\For{$t=1,\ldots,t_{\max}$}
  \State $\bm q\gets\nabla F\left(\widehat{\bm u}^{(t-1)}\right)$
  \State $\bm\mu\gets\beta_1\bm\mu+(1-\beta_1)\bm q$
  \State $\bm s\gets\bm s-(1-\beta_2)\operatorname{sign}\left(\bm s-\bm q^2\right)\bm q^2$
  \State $D\gets\operatorname{diag}\left(\sqrt{\bm s}+\epsilon\right)$
  \State $\widehat{\bm u}^{(t)}\gets\Pi^D_{\mathcal{S}}\left(\widehat{\bm u}^{(t-1)}-\eta D^{-1}\bm\mu/\left(1-\beta_1^t\right)\right)$
  \If{$r\left(\widehat{\bm u}^{(t)}\right)<\mathrm{tol}$}
    \State \textbf{break}
  \EndIf
\EndFor
\State \Return $\widehat{\bm u}^{(t)}$, $\rho=\exp F\left(\widehat{\bm u}^{(t)}\right)$
\end{algorithmic}
\end{algorithm}

\begin{proposition}\label{t-yogi-limit}
Let $\widehat{\bm u}^{(t)}$, $\bm\mu$, $\bm s$ be the iterates of
Algorithm~\ref{t-alg-yogi} with $t_{\max}=\infty$ and $\mathrm{tol}=0$, and
$G$ as in~\eqref{t-G}.
\begin{enumerate}
\item[(a)] Every iterate lies in $\mathcal{S}$, and
$0\le s_l\le S:=\max\left(\left\|\bm s^{(0)}\right\|_\infty,(2-\beta_2)G^2\right)$
for every entry $l$ and every $t$. Hence
$\epsilon I\le D\le\left(\sqrt{S}+\epsilon\right)I$.
\item[(b)] If $\widehat{\bm u}^{(t)}\to\widehat{\bm u}^*$, then
$\widehat{\bm u}^*$ satisfies the Karush--Kuhn--Tucker conditions
of~\eqref{M-three-program}.
\end{enumerate}
\end{proposition}
\begin{proof}
(a) The iterates lie in $\mathcal{S}$ because they are values of
$\Pi^D_{\mathcal{S}}$. Fix an entry $l$, write $s$ for $s_l$ before the update
and $q^2\le G^2$ for $q_l^2$. If $s>q^2$, the update gives
$s-(1-\beta_2)q^2$, which lies in $\left(\beta_2q^2,s\right)$ if $q\ne 0$ and
equals $s$ if $q=0$. If $s<q^2$, it gives $s+(1-\beta_2)q^2\in\left[0,(2-\beta_2)q^2\right]$
when $s\ge 0$. If $s=q^2$, $s$ is unchanged. By induction from
$\bm s^{(0)}=\bm q^2\ge 0$, each entry stays non-negative, and it either does
not increase or ends at most $(2-\beta_2)G^2$; hence it never exceeds $S$.

(b) By continuity of $\nabla F$,
$\bm q^{(t)}=\nabla F\left(\widehat{\bm u}^{(t-1)}\right)\to\nabla F\left(\widehat{\bm u}^*\right)$.
The moment
\begin{equation*}
\bm\mu^{(t)}=(1-\beta_1)\sum_{\tau=1}^t\beta_1^{t-\tau}\bm q^{(\tau)}
\end{equation*}
is a weighted mean with weights summing to $1-\beta_1^t$, so
$\bm\mu^{(t)}/\left(1-\beta_1^t\right)\to\nabla F\left(\widehat{\bm u}^*\right)$.
By (a) the matrices $D^{(t)}$ lie in a compact set of positive diagonal
matrices, so a subsequence converges to some $D^*\ge\epsilon I$. Along this
subsequence, the continuity of Proposition~\ref{t-proj-prop}(c) in the update
of $\widehat{\bm u}^{(t)}$ gives
\begin{equation*}
\widehat{\bm u}^*=\Pi^{D^*}_{\mathcal{S}}\left(\widehat{\bm u}^*-\eta D^{*-1}\nabla F\left(\widehat{\bm u}^*\right)\right),
\end{equation*}
and Proposition~\ref{t-fixed} applies.
\end{proof}

Proposition~\ref{t-yogi-limit} does not assert that the iterates converge, nor
that $F$ decreases along them; the analysis of~\cite{2018Zaheer_Yogi} is for
unconstrained problems with stochastic gradients and does not cover the
projection. Like the joint method, the method is local, and on the non-convex
problem~\eqref{M-three-program} its limit, when there is one, depends on the
start. Whether its steps, which treat all three blocks at once, avoid the
non-global endpoints of Algorithm~\ref{t-alg-gs} in Section~\ref{t-results} is
a question for the computations below.

\subsection{Results}\label{t-yogi-results}

In all runs $\eta=10^{-2}$, $\beta_1=0.9$, $\beta_2=0.999$,
$\epsilon=10^{-8}$, $\mathrm{tol}=10^{-10}$ and $t_{\max}=2\cdot 10^5$; no run
reached $t_{\max}$. With $\eta=10^{-3}$ the run from the barycenter of the
weekly instance had $r=1.3\cdot 10^{-7}$ after $5000$ iterations, against
convergence in $652$ iterations with $\eta=10^{-2}$.

\paragraph{The weekly instance.} On the instance of Section~\ref{t-results},
from the $49$ starts of Algorithm~\ref{t-alg-starts},
Algorithm~\ref{t-alg-yogi} ends at $\rho=1.661826042009$ from every start, the
values agreeing to $10^{-12}$; this is the smallest value of
Table~\ref{t-tab}. At every endpoint $r<10^{-10}$, and Algorithm~\ref{t-alg-kkt}
gives spread $\le 3.5\cdot 10^{-8}$ and excess $\ge 1.0\cdot 10^{-3}$, so the
check is conclusive, unlike for the joint method. In particular the $16$
starts with dose $3$ at a vertex, from which Algorithm~\ref{t-alg-gs} stops
at $\rho$ between $1.77$ and $1.95$, all reach the smallest value.
Table~\ref{t-tab-yogi} gives the iterations and the times.

\begin{table}[ht]
\centering
\small
\begin{tabular}{lcccc}
\hline
starts & barycenter & dose $1$ at a vertex & dose $2$ at a vertex & dose $3$ at a vertex\\
 & ($1$) & ($16$) & ($16$) & ($16$)\\
\hline
iterations & $652$ & $3159$--$15460$ & $12699$--$166447$ & $1118$--$3354$\\
time & $0.5$ s & $90$ s & $822$ s & $27$ s\\
\hline
\end{tabular}
\caption{Algorithm~\ref{t-alg-yogi} on the weekly instance from the $49$ starts
of Algorithm~\ref{t-alg-starts}; times are totals over the starts of the
column. Every run ends at $\rho=1.661826042009$.}\label{t-tab-yogi}
\end{table}

The slow runs are the $11$ starts with dose $2$ at a vertex in locations $2$,
$3$ and $4$, which took between $20546$ and $166447$ iterations; with
$t_{\max}=2\cdot 10^4$ eight of them stop at
$r$ between $4\cdot 10^{-6}$ and $1.2\cdot 10^{-4}$, with $\rho$ at most
$3.3\cdot 10^{-7}$ relative above the final value. This is the slow final
phase of a first-order method, not a different limit.

\paragraph{The $36$ instances.} On the instances of Section~\ref{t-instances},
from the barycenter, Algorithm~\ref{t-alg-yogi} takes between $545$ and
$10897$ iterations, $71$~s for all $36$, and ends with $r<10^{-10}$ and spread
$\le 4.0\cdot 10^{-8}$ in every instance. Its $\rho$ is within $10^{-7}$ of
$\rho_{\mathrm{ref}}$, or below it, in $31$ instances: $14$ of the $18$ weekly
and $17$ of the $18$ daily. Table~\ref{t-tab-yogi-inst} lists the others and
the two instances in which it ends below $\rho_{\mathrm{ref}}$.

\begin{table}[ht]
\centering
\small
\begin{tabular}{lllccc}
\hline
budget & $(b_{2i},b_{3i})/b_{1i}$ & $\sigma_2,\sigma_3$ & $\rho_{\mathrm{ref}}$ &
$\rho$, Alg.~\ref{t-alg-yogi} & relative excess\\
\hline
weekly & $(\tfrac14,\tfrac14)$       & seed $2$ & $1.7117509539$ & $1.7095128048$ & $-1.3\cdot 10^{-3}$\\
weekly & $(\tfrac18,\tfrac12)$       & $P$      & $1.3088277638$ & $1.3084969664$ & $-2.5\cdot 10^{-4}$\\
\hline
weekly & $(\tfrac1{16},\tfrac1{16})$ & $P$      & $2.5031932862$ & $2.5334636776$ & $1.2\cdot 10^{-2}$\\
weekly & $(\tfrac1{16},\tfrac1{16})$ & seed $2$ & $2.4252749378$ & $2.4383599028$ & $5.4\cdot 10^{-3}$\\
weekly & $(\tfrac18,\tfrac12)$       & seed $1$ & $1.2363297040$ & $1.2376346772$ & $1.1\cdot 10^{-3}$\\
weekly & $(\tfrac18,\tfrac12)$       & seed $2$ & $1.3770469379$ & $1.3819571439$ & $3.6\cdot 10^{-3}$\\
daily  & $(\tfrac1{16},\tfrac1{16})$ & $P$      & $4.0391397841$ & $4.0402962781$ & $2.9\cdot 10^{-4}$\\
\hline
\end{tabular}
\caption{Algorithm~\ref{t-alg-yogi} from the barycenter: the instances in which
it ends below $\rho_{\mathrm{ref}}$ (top) and those in which it is not
successful (bottom). The relative excess is
$(\rho-\rho_{\mathrm{ref}})/\rho_{\mathrm{ref}}$.}\label{t-tab-yogi-inst}
\end{table}

In two weekly instances Algorithm~\ref{t-alg-yogi} from the barycenter ends
below $\rho_{\mathrm{ref}}$, by $0.13$ and $0.025$ per cent, at points
satisfying the Karush--Kuhn--Tucker conditions to $r<10^{-10}$. The value
$\rho_{\mathrm{ref}}$ of Section~\ref{t-instances}, the best of $69$ runs of
the joint method, is therefore not the global minimum there. In both instances
the joint method and Algorithm~\ref{t-alg-gs} from the barycenter had ended at
$\rho_{\mathrm{ref}}$; measured against the smaller value they are not
successful, which lowers the weekly counts of Table~\ref{t-tab-inst} from $17$
to $15$ for the joint method and from $16$ to $14$ for
Algorithm~\ref{t-alg-gs}. In the two weekly instances with
$b_{2i}=b_{3i}=b_{1i}/16$ Algorithm~\ref{t-alg-yogi} ends at the same value as
Algorithm~\ref{t-alg-gs}, and in the second of them also as the joint method.
Of the daily instances, it misses $\rho_{\mathrm{ref}}$ only in the one in which
the two other methods also miss it, by $2.9\cdot 10^{-4}$ against
$1.1\cdot 10^{-4}$ for each of them.

\subsection{Convergence and allocations}\label{t-yogi-figures}

Figure~\ref{t-fig-conv} shows, on the weekly instance of
Section~\ref{t-results}, the relative excess $(\rho-\rho^*)/\rho^*$ over the
smallest value $\rho^*=1.661826042009$ and the stationarity measure $r$
of~\eqref{t-r} along Algorithm~\ref{t-alg-yogi} from three starts: the
barycenter, the start $(k,i,j)=(3,1,4)$, from which Algorithm~\ref{t-alg-gs}
stops at $\rho=1.9509854920$, and the start $(2,4,1)$, the slowest of
Table~\ref{t-tab-yogi}. The relative excess falls below $10^{-8}$ after $174$
iterations from the barycenter and after $340$ from the start $(3,1,4)$, from
which Algorithm~\ref{t-alg-gs} stalls after two sweeps at the relative excess
$0.174$. From the start $(2,4,1)$ it falls below $10^{-8}$ only after $37290$
iterations, and $r$ decreases slowly over most of the run, which is the slow
final phase of Section~\ref{t-yogi-results}. Algorithm~\ref{t-alg-gs} from
the barycenter reaches the relative excess $1.2\cdot 10^{-10}$ after $5$
sweeps, each of which solves three block problems.

\pgfplotsset{conv panel/.style={
  width=0.345\textwidth, height=0.34\textwidth,
  ymode=log, ymin=1e-16, ymax=2,
  ymajorgrids, grid style={gray!25},
  axis lines*=left, title style={font=\small},
  tick label style={font=\scriptsize}, label style={font=\small},
  legend cell align=left,
  legend style={draw=none, fill=white, fill opacity=0.8, text opacity=1, font=\scriptsize},
  every axis plot/.append style={line width=1.0pt}}}
\begin{figure}[ht]
\centering
\begin{tikzpicture}
\begin{groupplot}[group style={group size=3 by 1, horizontal sep=1.1cm}, conv panel]
\nextgroupplot[title={Algorithm~\ref{t-alg-yogi}: relative excess}, xmode=log,
  xlabel={iteration}, ylabel={$(\rho-\rho^*)/\rho^*$}, legend pos=south west]
\addplot[color=locNE] table[x=it, y=excess]{fig_yogi_hist_bary.dat};
\addplot[color=locMW] table[x=it, y=excess]{fig_yogi_hist_k3.dat};
\addplot[color=locS]  table[x=it, y=excess]{fig_yogi_hist_k2.dat};
\legend{barycenter, {$(3,1,4)$}, {$(2,4,1)$}}
\nextgroupplot[title={Algorithm~\ref{t-alg-yogi}: $r$}, xmode=log,
  xlabel={iteration}]
\addplot[color=locNE] table[x=it, y=r]{fig_yogi_hist_bary.dat};
\addplot[color=locMW] table[x=it, y=r]{fig_yogi_hist_k3.dat};
\addplot[color=locS]  table[x=it, y=r]{fig_yogi_hist_k2.dat};
\nextgroupplot[title={Algorithm~\ref{t-alg-gs}: relative excess},
  xlabel={sweep}, xmin=0, xmax=16, legend pos=south east]
\addplot[color=locNE, mark=*, mark size=1.5pt] table[x=sweep, y=excess]{fig_gs_hist_bary.dat};
\addplot[color=locMW, mark=square*, mark size=1.5pt] table[x=sweep, y=excess]{fig_gs_hist_k3.dat};
\legend{barycenter, {$(3,1,4)$}}
\end{groupplot}
\end{tikzpicture}
\caption{The weekly instance of Section~\ref{t-results}. Left and middle:
Algorithm~\ref{t-alg-yogi} from the barycenter and from the starts
$(k,i,j)=(3,1,4)$ and $(2,4,1)$ of Algorithm~\ref{t-alg-starts}, the relative
excess over $\rho^*=1.661826042009$ and the stationarity measure
$r$ of~\eqref{t-r} per iteration; the iteration count is on a logarithmic scale,
and the start, iteration $0$, is not shown. Right: Algorithm~\ref{t-alg-gs}
from the barycenter and from $(3,1,4)$, the relative excess per sweep. Excesses
below $10^{-16}$, rounding, are shown at $10^{-16}$.}\label{t-fig-conv}
\end{figure}

Figures~\ref{t-fig-alloc-s2} and~\ref{t-fig-alloc-p} show the allocations
$\chi_{k,ij}$ at the endpoints of the joint method and of
Algorithm~\ref{t-alg-yogi}, both from the barycenter, in the two instances of
Table~\ref{t-tab-yogi-inst} in which Algorithm~\ref{t-alg-yogi} ends below
$\rho_{\mathrm{ref}}$. The joint method ends at $\rho_{\mathrm{ref}}$ in both.
The two allocations differ mainly in one location. In the first instance the
endpoint of the joint method gives the group $60+$ of the South
$\chi_{2}=0.0162$ and $\chi_{3}=0.0211$ per day, and that of
Algorithm~\ref{t-alg-yogi} gives this group no second or third doses and their
budget to the groups $0$--$59$ of the South. In the second instance the same
holds for the group $60+$ of the West, with $\chi_{2}=0.0098$ and
$\chi_{3}=0.0181$ per day at the endpoint of the joint method. In the other
three locations the entries above $10^{-3}$ per day agree to $2.2$ per cent;
the smaller ones, first doses to the group $60+$, differ by up to
$3\cdot 10^{-4}$ per day. As for the single dose,
where the group $60+$ receives no doses at the minimizer, the smaller value
moves doses away from this group.

\pgfplotsset{alloc panel/.style={
  width=0.46\textwidth, height=0.26\textwidth,
  xtick={0,1,2,3}, xticklabels={0--19,20--39,40--59,60+},
  ymin=0, ymajorgrids, grid style={gray!25},
  axis lines*=left, title style={font=\small},
  tick label style={font=\scriptsize}, label style={font=\small},
  legend cell align=left,
  legend style={draw=none, fill=none, font=\scriptsize},
  every axis plot/.append style={line width=1.0pt, mark size=1.8pt}}}
% #1: file stem, #2: rho of the joint method, #3: rho of Yogi
\newcommand{\allocfigure}[3]{%
\begin{tikzpicture}
\begin{groupplot}[group style={group size=2 by 3, horizontal sep=1.2cm,
    vertical sep=1.3cm, xticklabels at=edge bottom}, alloc panel]
\nextgroupplot[title={joint method, $\rho=#2$}, ymax=6, ylabel={$\chi_1$}, legend pos=south west]
\addplot[color=locNE, mark=*]         table[x=age, y=k1NE]{fig_alloc_#1_joint.dat};
\addplot[color=locMW, mark=square*]   table[x=age, y=k1MW]{fig_alloc_#1_joint.dat};
\addplot[color=locS,  mark=triangle*] table[x=age, y=k1S]{fig_alloc_#1_joint.dat};
\addplot[color=locW,  mark=diamond*]  table[x=age, y=k1W]{fig_alloc_#1_joint.dat};
\legend{Northeast, Midwest, South, West}
\nextgroupplot[title={Algorithm~\ref{t-alg-yogi}, $\rho=#3$}, ymax=6]
\addplot[color=locNE, mark=*]         table[x=age, y=k1NE]{fig_alloc_#1_yogi.dat};
\addplot[color=locMW, mark=square*]   table[x=age, y=k1MW]{fig_alloc_#1_yogi.dat};
\addplot[color=locS,  mark=triangle*] table[x=age, y=k1S]{fig_alloc_#1_yogi.dat};
\addplot[color=locW,  mark=diamond*]  table[x=age, y=k1W]{fig_alloc_#1_yogi.dat};
\nextgroupplot[ymax=60, ylabel={$\chi_2$}]
\addplot[color=locNE, mark=*]         table[x=age, y=k2NE]{fig_alloc_#1_joint.dat};
\addplot[color=locMW, mark=square*]   table[x=age, y=k2MW]{fig_alloc_#1_joint.dat};
\addplot[color=locS,  mark=triangle*] table[x=age, y=k2S]{fig_alloc_#1_joint.dat};
\addplot[color=locW,  mark=diamond*]  table[x=age, y=k2W]{fig_alloc_#1_joint.dat};
\nextgroupplot[ymax=60]
\addplot[color=locNE, mark=*]         table[x=age, y=k2NE]{fig_alloc_#1_yogi.dat};
\addplot[color=locMW, mark=square*]   table[x=age, y=k2MW]{fig_alloc_#1_yogi.dat};
\addplot[color=locS,  mark=triangle*] table[x=age, y=k2S]{fig_alloc_#1_yogi.dat};
\addplot[color=locW,  mark=diamond*]  table[x=age, y=k2W]{fig_alloc_#1_yogi.dat};
\nextgroupplot[ymax=100, ylabel={$\chi_3$}, xlabel={age group (years)}]
\addplot[color=locNE, mark=*]         table[x=age, y=k3NE]{fig_alloc_#1_joint.dat};
\addplot[color=locMW, mark=square*]   table[x=age, y=k3MW]{fig_alloc_#1_joint.dat};
\addplot[color=locS,  mark=triangle*] table[x=age, y=k3S]{fig_alloc_#1_joint.dat};
\addplot[color=locW,  mark=diamond*]  table[x=age, y=k3W]{fig_alloc_#1_joint.dat};
\nextgroupplot[ymax=100, xlabel={age group (years)}]
\addplot[color=locNE, mark=*]         table[x=age, y=k3NE]{fig_alloc_#1_yogi.dat};
\addplot[color=locMW, mark=square*]   table[x=age, y=k3MW]{fig_alloc_#1_yogi.dat};
\addplot[color=locS,  mark=triangle*] table[x=age, y=k3S]{fig_alloc_#1_yogi.dat};
\addplot[color=locW,  mark=diamond*]  table[x=age, y=k3W]{fig_alloc_#1_yogi.dat};
\end{groupplot}
\end{tikzpicture}}

\begin{figure}[ht]
\centering
\allocfigure{s2}{1.7117509539}{1.7095128048}
\caption{The weekly instance with $b_{2i}=b_{3i}=b_{1i}/4$ and $\sigma_2$,
$\sigma_3$ of seed $2$: the allocations $\chi_1$, $\chi_2$, $\chi_3$ (rows),
in units of $10^{-3}$ per day, at the endpoints of the joint method (left) and
of Algorithm~\ref{t-alg-yogi} (right), both from the barycenter. One line per
location; the same scale in each row.}\label{t-fig-alloc-s2}
\end{figure}

\begin{figure}[ht]
\centering
\allocfigure{p}{1.3088277638}{1.3084969664}
\caption{As Figure~\ref{t-fig-alloc-s2}, for the weekly instance with
$b_{2i}=b_{1i}/8$, $b_{3i}=b_{1i}/2$ and $\sigma_2=\sigma_3=P$.}\label{t-fig-alloc-p}
\end{figure}

\section{Conclusions}\label{t-conclusions}

\begin{enumerate}
\item The three-dose problem~\eqref{M-three-program} is not convex for the
weekly instance either: it has local minimizers with $\rho$ between $1.77$ and
$1.95$ besides the smallest value $\rho=1.6618260420$ found. Theorems~\ref{M-global}
and~\ref{M-barrier} have no analogue here, and a computed point satisfying the
Karush--Kuhn--Tucker conditions need not be a global minimizer.
\item Of the three blocks, the first is strictly convex
(Proposition~\ref{t-block1}) and the third is convex
(Proposition~\ref{t-block3}); the second is neither, and its dose $\chi_2$ couples
to $\chi_3$ through the product $\pi_{1,p}\pi_{2,p}\pi_{3,p}$ in~\eqref{M-g-def}.
\item Block Gauss--Seidel, Algorithm~\ref{t-alg-gs}, decreases $\rho$ monotonically
and reaches the smallest value from the barycenter and from the $32$ vertex
starts with dose $1$ or $2$ at a vertex, at points that pass the
Karush--Kuhn--Tucker check. From the $16$ starts with dose $3$ at a vertex it
stops at non-global local minimizers: when
the second and third doses of a location are concentrated in one age group,
moving either of them alone does not decrease $\rho$.
\item On the weekly instance of Section~\ref{t-results} the joint method
reached the smallest value from all $49$ starts, in $25$~s in total against
$328$~s for Algorithm~\ref{t-alg-gs}.
\item Over $36$ instances (Section~\ref{t-instances}) the barycenter is a good
start in most cases but not in all: the joint method from it reaches the best
value found in $33$ instances, Algorithm~\ref{t-alg-gs} in $28$, and the better
of the two in $34$. The misses are up to $0.54$ and $1.2$ per cent in $\rho$,
and the hard instances are those with small second-dose budgets and daily
first-dose budgets.
\item In practice the joint method should be run from several starts, the
barycenter, the starts near the vertices and random points, and the smallest
value kept; neither method, from any single start, is guaranteed to find the
global minimum. These tests use second and third dose data chosen for lack of
data, $\sigma_2$ and $\sigma_3$ built from the numbers of partially protected
individuals.
\end{enumerate}

The computations are in the following files:
\begin{itemize}
\item \texttt{gs\_three\_dose.py}: Algorithms~\ref{t-alg-gs}, \ref{t-alg-kkt}
and~\ref{t-alg-starts}, and Table~\ref{t-tab};
\item \texttt{check\_block2.py}: the sampling for the second block;
\item \texttt{check\_gs\_endpoints.py}: the two endpoints and the convexity test;
\item \texttt{test\_barycenter\_three\_dose.py}: the $36$ instances of
Section~\ref{t-instances} and Tables~\ref{t-tab-inst} and~\ref{t-tab-fail};
\item \texttt{yogi\_three\_dose.py}: Algorithm~\ref{t-alg-yogi} from the $49$
starts, and Table~\ref{t-tab-yogi};
\item \texttt{yogi\_instances\_three\_dose.py}: Algorithm~\ref{t-alg-yogi} on
the $36$ instances, and Table~\ref{t-tab-yogi-inst};
\item \texttt{plot\_yogi\_three\_dose.py}: the data of
Figures~\ref{t-fig-conv}, \ref{t-fig-alloc-s2} and~\ref{t-fig-alloc-p}.
\end{itemize}

\bibliographystyle{plain}
\bibliography{bib_mb}

\end{document}
